\documentclass[../../main.tex]{subfiles}
\begin{document}
% SOURCE p=206; printed=191; object=ch06-dual-exercises-title; region=lower
\Needspace{12\baselineskip}
\subsection*{习题}
% SOURCE p=206; printed=191; object=ch06-ex17; number=17; region=lower
\begin{exercise}\label{exercise:ch06-test-sigma-support}
沿用\cref{theorem:ch06-dual-test-criterion}的记号，若 $\mu$ 半有限、$q<\infty$ 且 $M_q(g)<\infty$，则对每个 $\varepsilon>0$，$\{x:|g(x)|>\varepsilon\}$ 都有有限测度，从而 $S_g$ 为 $\sigma$ 有限集。
\end{exercise}
% SUPPLEMENT: id=v14-ch06-ex17-solution; parent=ch06-ex17; kind=solution+motivation+hypotheses; origin=independent-derivation
\begin{solution}
\textbf{不依赖待证对偶判别的证明（重构新增）。}令 $M=M_q(g)$、$A=\{|g|>\varepsilon\}$。对任意可测 $B\subset A$，若 $0<\mu(B)<\infty$，取
\begin{align*}
u=\mu(B)^{-\frac1p}\chi_B\overline{\operatorname{sgn}g},
\end{align*}
其中 $p=\infty$ 时 $\mu(B)^{-1/p}=1$。则 $\|u\|_p=1$。虽 $u$ 未必简单，但用受 $|u|$ 控制的简单函数逼近，再以 $\mu(B)^{-1/p}|g|\chi_B\in L^1$ 作控制函数，得到 $|\int ug|\leqslant M$。这只用原假设 $\chi_Bg\in L^1$ 和 $M$ 的定义。因此
\begin{align*}
M\geqslant\mu(B)^{-\frac1p}\int_B|g|
\geqslant\varepsilon\mu(B)^{\frac1q},
\qquad \mu(B)\leqslant(M/\varepsilon)^q.
\end{align*}
零测度的 $B$ 也满足最后的界。

为从有限子集推出整个 $A$ 有限，令
\begin{align*}
C=\sup\{\mu(B):B\subset A,\ B\text{ 可测},\ \mu(B)<\infty\}\leqslant(M/\varepsilon)^q.
\end{align*}
选有限测度 $B_n\subset A$，使 $\mu(B_n)\to C$；$C=0$ 时可全取空集。令 $F_n=\bigcup_{j=1}^nB_j$，则 $F_n$ 有限测度，故 $\mu(F_n)\leqslant C$。$F=\bigcup_nF_n$ 满足 $\mu(F)=C<\infty$。若 $\mu(A\setminus F)>0$，则该集合本身有限且正时可直接取它，无穷时由半有限性取正有限测度子集 $D$。于是 $F\cup D\subset A$ 有限测度且测度大于 $C$，矛盾。所以 $\mu(A\setminus F)=0$，$\mu(A)=C<\infty$。最后 $S_g=\bigcup_{k\geqslant1}\{|g|>1/k\}$ 为 $\sigma$ 有限集。
\end{solution}
\begin{remark}
\textbf{动机与条件说明（重构新增）。}有限测试的上界控制每个水平集所能容纳的有限质量；半有限性排除仍藏有不可探测的正质量。$q<\infty$ 使 $\mu(B)^{1/q}$ 随测度增长，给统一有限界；$q=\infty$ 时这个机制消失，如无限测度空间上的 $g=1$ 可有有限测试上界而正水平集无限。$\varepsilon>0$ 用于除法，并把 $g$ 与零分开；$M=0$ 仍包含在论证中。这里没有引用 $g\in L^q$，因而可以作为正文判别定理的独立前置证明。
\end{remark}
% END SUPPLEMENT: id=v14-ch06-ex17-solution
% SOURCE p=206; printed=191; object=ch06-ex18; number=18; region=bottom
\begin{exercise}\label{exercise:ch06-radon-nikodym-hilbert}
$L^2$ 的自对偶性由 Hilbert 空间理论（\cref{theorem:ch05-riesz-hilbert}）得出；利用这一事实，可以通过下面由 von Neumann 提出的论证证明 Lebesgue–Radon–Nikodym 定理。设 $\mu,\nu$ 是 $(X,\mathscr M)$ 上的有限正测度（$\sigma$ 有限情形可像\cref{section:ch03-lebesgue-radon-nikodym}中那样容易地推出），令 $\lambda=\mu+\nu$。
\begin{enumerate}
\item\label{item:ch06-rn-hilbert-representation} 映射 $f\mapsto\int f\,\mathrm d\nu$ 是 $L^2(\lambda)$ 上的有界线性泛函，所以对某个 $g\in L^2(\lambda)$，有 $\int f\,\mathrm d\nu=\int fg\,\mathrm d\lambda$。等价地，对所有 $f\in L^2(\lambda)$，
\begin{align*}
\int f(1-g)\,\mathrm d\nu=\int fg\,\mathrm d\mu.
\end{align*}
\item\label{item:ch06-rn-hilbert-bounds} $0\leqslant g\leqslant1$，$\lambda$ 几乎处处成立，因此可假设它处处成立。
\item\label{item:ch06-rn-hilbert-decomposition} 令 $A=\{x:g(x)<1\}$、$B=\{x:g(x)=1\}$，并令 $\nu_a(E)=\nu(A\cap E)$、$\nu_s(E)=\nu(B\cap E)$。则 $\nu_s\perp\mu$ 且 $\nu_a\ll\mu$；事实上，
\begin{align*}
\mathrm d\nu_a=g(1-g)^{-1}\chi_A\,\mathrm d\mu.
\end{align*}
\end{enumerate}
\end{exercise}
% SUPPLEMENT: id=v14-ch06-ex18-solution; parent=ch06-ex18; kind=solution+motivation+hypotheses; origin=independent-derivation
\begin{solution}
\textbf{三问及推广的独立证明（重构新增）。}
\begin{enumerate}
\item 因 $\nu\leqslant\lambda$，$\lambda$ 零集也是 $\nu$ 零集，泛函与代表无关。由 Schwarz 不等式，
\begin{align*}
\left|\int f\,\mathrm d\nu\right|
\leqslant\nu(X)^{\frac12}\left(\int|f|^2\,\mathrm d\nu\right)^{\frac12}
\leqslant\nu(X)^{\frac12}\|f\|_{L^2(\lambda)}.
\end{align*}
\cref{theorem:ch06-lp-completeness}保证 $L^2(\lambda)$ 是 Hilbert 空间；Hilbert 空间 Riesz 定理给 $h\in L^2(\lambda)$，使泛函等于 $\int f\overline h\,\mathrm d\lambda$。取 $g=\overline h$ 即得所需形式。$fg\in L^1(\lambda)$，从而也对 $\mu,\nu$ 可积，展开 $\lambda=\mu+\nu$ 并移项合法，得到所述等价等式。此处不使用本节的 $L^p$ 对偶表示定理，也未使用 Radon–Nikodym 定理。
\item 因 $\lambda(X)<\infty$，每个 $\chi_E$ 属于 $L^2(\lambda)$，且 $g\in L^1(\lambda)$。代入得
\begin{align*}
\int_Eg\,\mathrm d\lambda=\nu(E)\in[0,\lambda(E)]
\quad(E\in\mathscr M).
\end{align*}
所有可测集上 $\operatorname{Im}g$ 的积分为零，故其正、负水平集 $\{\operatorname{Im}g>1/k\}$、$\{\operatorname{Im}g<-1/k\}$ 均为零集：否则对应积分严格正或负。于是 $g$ 为实值，$\lambda$ 几乎处处成立。在该满测度集上，若 $\{g<-1/k\}$ 正测度，其积分为负，与 $\nu\geqslant0$ 矛盾；若 $\{g>1+1/k\}$ 正测度，则其上 $g$ 的积分严格大于该集的 $\lambda$ 测度，也矛盾。取可数并，得 $0\leqslant g\leqslant1$ 几乎处处。把可测例外零集上的 $g$ 改为零，便得到处处界，同时不改变第一问等式。
\item 代 $f=\chi_B$ 入第一问的等价式，左侧为零，右侧为 $\mu(B)$，所以 $\mu(B)=0$。$\nu_s$ 集中于 $B$，故 $\nu_s\perp\mu$。在 $A$ 上令 $h=g/(1-g)$，在 $B$ 上令 $h=0$；这是非负可测函数，不在 $B$ 上计算 $0\cdot\infty$。任取可测 $E$，令
\begin{align*}
u_k=\chi_{E\cap A}\min\left(k,\frac1{1-g}\right),
\end{align*}
并在 $B$ 上直接定义为零。$u_k$ 有界且 $\lambda$ 有限，故属于 $L^2(\lambda)$。代入第一问得到
\begin{align*}
\int_{E\cap A}\min\bigl(k(1-g),1\bigr)\,\mathrm d\nu
=\int_{E\cap A}g\min\left(k,\frac1{1-g}\right)\,\mathrm d\mu.
\end{align*}
两边的被积函数都随 $k$ 单调增加，分别趋向 $\chi_{E\cap A}$ 与 $h\chi_E$。由单调收敛，$\nu_a(E)=\int_Eh\,\mathrm d\mu$，因而 $\nu_a\ll\mu$。$A,B$ 为可测分割，所以 $\nu=\nu_a+\nu_s$。
\end{enumerate}
为补全题中提及的 $\sigma$ 有限推广，若 $\mu,\nu$ 均 $\sigma$ 有限，分别取有限测度穷竭，将它们两两相交，再化为可数可测分割 $X=\bigsqcup_jD_j$，使每个 $D_j$ 对两测度均有限。逐块应用上面的有限情形，得 $\nu|_{D_j}=h_j\mu|_{D_j}+s_j$，其中 $s_j$ 集中于 $D_j$ 中的 $\mu$ 零集 $N_j$。令 $h=h_j$ 于 $D_j$，并令 $s(E)=\sum_js_j(E\cap D_j)$。非负函数的单调收敛和测度的可数可加性给 $\nu=h\mu+s$；$s$ 集中于零集 $\bigcup_jN_j$，故 $s\perp\mu$。

唯一性也可直接证明。若 $\nu=h\mu+s=k\mu+t$，其中 $s,t\perp\mu$，取一个共同的 $\mu$ 零集 $N$，使 $s,t$ 都集中于 $N$。在 $N^c$ 上，两个密度对每个可测集的积分相同。在上述共同有限分块上，分别测试 $\{h\geqslant k+1/n\}$ 和 $\{k\geqslant h+1/n\}$，可得这些集合均为 $\mu$ 零集，从而 $h=k$ 几乎处处。密度在 $N$ 上的积分为零，故 $s(E)=\nu(E\cap N)=t(E)$。这也证明绝对连续情形的奇异部分必为零。符号测度版本再对 Jordan 分解的两部分分别应用此正测度结果。
\end{solution}
\begin{remark}
\textbf{动机与逐条件说明（重构新增）。}先用共同支配测度 $\lambda=\mu+\nu$ 把积分泛函放进 Hilbert 空间；表示元 $g$ 衡量 $\nu$ 相对总量的比例。$g=1$ 的区域没有 $\mu$ 质量，因而产生奇异部分；$g<1$ 的区域可除以剩余比例，得到相对 $\mu$ 的密度。正性给 $0\leqslant g\leqslant1$ 和单调收敛；有限性让示性函数及截断倒数属于 $L^2$。不假设 $\nu\ll\mu$，所以不能预先删掉 $B$；直接使用未截断的 $(1-g)^{-1}$ 也不能保证它属于测试空间。$\sigma$ 有限推广需要共同的可数有限分块，使可数零集并仍零；零测度情形同样成立。
\end{remark}
% END SUPPLEMENT: id=v14-ch06-ex18-solution
% SOURCE p=207; printed=192; object=ch06-ex19; number=19; region=top
\begin{exercise}\label{exercise:ch06-linfty-dual-cluster}
以 $\phi_n(f)=n^{-1}\sum_{j=1}^nf(j)$ 定义 $\phi_n\in(\ell^\infty)^*$。则序列 $\{\phi_n\}$ 有弱星聚点 $\phi$，且 $\phi$ 是 $(\ell^\infty)^*$ 中不能由某个 $\ell^1$ 元素产生的元素。
\end{exercise}
% SUPPLEMENT: id=v14-ch06-ex19-solution; parent=ch06-ex19; kind=solution+motivation+hypotheses; origin=independent-derivation
\begin{solution}
\textbf{紧性与坐标测试（重构新增）。}有 $|\phi_n(f)|\leqslant\|f\|_\infty$，而在常数序列 $\boldsymbol1$ 上 $\phi_n(\boldsymbol1)=1$，所以 $\|\phi_n\|=1$。由\cref{theorem:ch05-alaoglu}，$(\ell^\infty)^*$ 的闭单位球 $K$ 弱星紧。对每个 $N$，令 $K_N$ 为 $\{\phi_n:n\geqslant N\}$ 在 $K$ 中的弱星闭包。这是非空闭集的递减列，有有限交性质，故存在 $\phi\in\bigcap_NK_N$。每个 $\phi$ 的弱星邻域都与任意尾部相交，因此它确为该序列的弱星聚点。

点评价 $\psi\mapsto\psi(f)$ 在弱星拓扑下连续。所有 $\phi_n(\boldsymbol1)=1$，故 $\phi(\boldsymbol1)=1$。对固定坐标向量 $e_j$，$\phi_n(e_j)$ 在 $n\geqslant j$ 时等于 $1/n$，故趋零；给定 $\varepsilon>0$，充分后尾部包含于闭集 $\{\psi:|\psi(e_j)|\leqslant\varepsilon\}$，因而 $|\phi(e_j)|\leqslant\varepsilon$，所以 $\phi(e_j)=0$。若某个 $a\in\ell^1$ 给 $\phi(f)=\sum_ja_jf(j)$，测试每个 $e_j$ 得 $a_j=0$；于是 $\phi(\boldsymbol1)=0$，矛盾。
\end{solution}
\begin{remark}
\textbf{动机与条件说明（重构新增）。}平均把每个固定坐标的权重逐渐稀释，却始终保留常数序列的值；聚点因而不能由绝对可和的固定坐标权重表示。归一化 $1/n$ 保证范数统一为一，Alaoglu 才提供弱星紧性。这里取对预对偶 $\ell^\infty$ 的弱星拓扑；紧性给聚点或收敛子网，并未保证存在收敛子序列。$\ell^1$ 假设确保所说的级数表示对每个有界序列连续且有意义。
\end{remark}
% END SUPPLEMENT: id=v14-ch06-ex19-solution
% SOURCE p=207; printed=192; object=ch06-ex20; number=20; region=upper
\begin{exercise}\label{exercise:ch06-ae-bounded-weak}
设 $\sup_n\|f_n\|_p<\infty$，且 $f_n\to f$ 几乎处处。
\begin{enumerate}
\item\label{item:ch06-ae-bounded-weak-interior} 若 $1<p<\infty$，则 $f_n\to f$ 于 $L^p$ 的弱拓扑。（给定 $g\in L^q$，其中 $q$ 与 $p$ 共轭，及 $\varepsilon>0$，存在：(i) $\delta>0$，使 $\mu(E)<\delta$ 时 $\int_E|g|^q<\varepsilon$；(ii) $A\subset X$，使 $\mu(A)<\infty$ 且 $\int_{X\setminus A}|g|^q<\varepsilon$；(iii) $B\subset A$，使 $\mu(A\setminus B)<\delta$，且 $f_n\to f$ 在 $B$ 上一致成立。）
\item\label{item:ch06-ae-bounded-weak-endpoints} 第一问在 $p=1$ 时一般不成立；在 $L^1(\mathbb R,m)$ 与 $\ell^1$ 中找反例。不过，若 $\mu$ 为 $\sigma$ 有限测度，并把弱收敛换成弱星收敛，则结论对 $p=\infty$ 成立。
\end{enumerate}
\end{exercise}
% SUPPLEMENT: id=v14-ch06-ex20-solution; parent=ch06-ex20; kind=solution+motivation+hypotheses; origin=independent-derivation
\begin{solution}
\textbf{三部分估计与两个端点（重构新增）。}
\begin{enumerate}
\item 令 $C=\sup_n\|f_n\|_p$。Fatou 引理给 $\|f\|_p\leqslant C$，因此 $h_n=f_n-f$ 满足 $\|h_n\|_p\leqslant2C$。若 $C=0$，所有 $h_n=0$ 几乎处处，结论直接成立。以下设 $C>0$，固定 $g\in L^q$ 与目标误差 $\eta>0$。$q<\infty$，所以 $|g|^q\in L^1$；其积分对小测度集合绝对连续，且可在有限测度集合上截断。具体地，前者由先截掉 $|g|^q$ 的高值尾部再控制有界部分得到；后者可取 $A_k=\{|g|>1/k\}$，其测度有限，而 $\int_{A_k^c}|g|^q\to0$。故可取 $\delta>0$ 与有限测度 $A$，使
\begin{align*}
\mu(E)<\delta\ \Longrightarrow\ \|g\chi_E\|_q<\frac{\eta}{6C},
\qquad \|g\chi_{A^c}\|_q<\frac{\eta}{6C}.
\end{align*}
由\cref{theorem:ch02-egoroff}，存在可测 $B\subset A$，$\mu(A\setminus B)<\delta$，且 $h_n\to0$ 在 $B$ 上一致。由于 $\int_A|g|\leqslant\mu(A)^{1/p}\|g\|_q<\infty$，充分大的 $n$ 满足 $\sup_B|h_n|\int_A|g|<\eta/3$；若此积分为零，该项直接为零。分三部分并用 Hölder 不等式，
\begin{align*}
\left|\int h_ng\right|
&\leqslant\|h_n\|_p\|g\chi_{A^c}\|_q
+\|h_n\|_p\|g\chi_{A\setminus B}\|_q
+\sup_B|h_n|\int_B|g|<\eta.
\end{align*}
对每个 $g\in L^q$ 的配对均趋零，\cref{theorem:ch06-lp-duality}遂给弱收敛。
\item 在 $L^1(\mathbb R,m)$ 取 $f_n=\chi_{(n,n+1)}$，则范数恒为一且逐点趋零，但连续泛函 $f\mapsto\int f$ 在其上恒为一，故不弱收敛到零。在 $\ell^1$ 取 $f_n=e_n$，同样范数一、逐坐标趋零，但坐标求和泛函恒为一。

当 $p=\infty$ 时，令 $C=\sup_n\|f_n\|_\infty$。删去可数个本质界例外零集以及收敛例外零集，可得同时 $|f_n|\leqslant C$、$|f|\leqslant C$，故 $f\in L^\infty$。对每个 $g\in L^1$，有 $|(f_n-f)g|\leqslant2C|g|\in L^1$，且乘积几乎处处趋零。控制收敛给 $\int(f_n-f)g\to0$。$\mu$ 为 $\sigma$ 有限测度时，$L^\infty=(L^1)^*$，上述收敛正是弱星收敛。
\end{enumerate}
\end{solution}
\begin{remark}
\textbf{动机与逐条件说明（重构新增）。}弱收敛只检验与固定权重的积分；范数有界防止函数在权重尾部及小例外集上的贡献失控，几乎处处收敛经 Egoroff 定理控制余下有限区域。内部指数使共轭 $q$ 有限，权重的 $q$ 次积分才有这两种截断性质。$p=1$ 的常数测试权重无法衰减，两个反例中质量逃向远处或高坐标。$p=\infty$ 的统一逐点界可直接使用控制收敛；其积分配对收敛不需 $\sigma$ 有限性，但把整个 $L^\infty$ 识别为 $(L^1)^*$ 时，本题采用了该充分条件。此处没有把弱星收敛提升为弱收敛，也不要求测度空间完备。
\end{remark}
% END SUPPLEMENT: id=v14-ch06-ex20-solution
% SOURCE p=207; printed=192; object=ch06-ex21; number=21; region=middle
\begin{exercise}\label{exercise:ch06-sequence-weak-convergence}
若 $1<p<\infty$，则 $f_n\to f$ 于 $\ell^p(A)$ 的弱拓扑，当且仅当 $\sup_n\|f_n\|_p<\infty$ 且 $f_n\to f$ 逐点成立。
\end{exercise}
% SUPPLEMENT: id=v14-ch06-ex21-solution; parent=ch06-ex21; kind=solution+motivation+hypotheses; origin=independent-derivation
\begin{solution}
\textbf{两方向完整证明（重构新增）。}若 $f_n\to f$ 弱收敛，则\subcref{exercise:ch05-weak-convergent-bounded}{item:ch05-weak-and-star-sequences-bounded}给范数有界。对每个 $a\in A$，坐标评价 $u\mapsto u(a)$ 线性且 $|u(a)|\leqslant\|u\|_p$，所以连续；弱收敛便给 $f_n(a)\to f(a)$。

反之，设 $C=\sup_n\|f_n\|_p<\infty$ 且逐点极限为 $f$。对任意有限 $F\subset A$，有限求和可过极限，故
\begin{align*}
\sum_{a\in F}|f(a)|^p=\lim_n\sum_{a\in F}|f_n(a)|^p\leqslant C^p.
\end{align*}
取所有有限 $F$ 的上确界得 $f\in\ell^p(A)$。计数测度下逐点收敛就是几乎处处收敛，所以\subcref{exercise:ch06-ae-bounded-weak}{item:ch06-ae-bounded-weak-interior}给弱收敛。该结论适用于任意测度空间，故即使 $A$ 不可数也可应用。
\end{solution}
\begin{remark}
\textbf{动机与条件说明（重构新增）。}在序列空间中，单个坐标的质量不能隐藏在非空零集中，因此弱收敛可以用“逐坐标收敛加统一范数控制”刻画。必要方向使用的是弱收敛的范数有界性，不能只由每个坐标有界推得范数有界。内部指数用于充分方向；$p=1$ 的 $e_n$ 已给反例。$A$ 无需可数：有限子和定义范数，有限指数函数的非零坐标自动至多可数。
\end{remark}
% END SUPPLEMENT: id=v14-ch06-ex21-solution
% SOURCE p=207; printed=192; object=ch06-ex22; number=22; region=middle
\begin{exercise}\label{exercise:ch06-weak-ae-measure-examples}
令 $X=[0,1]$，配 Lebesgue 测度。
\begin{enumerate}
\item\label{item:ch06-cosine-weak-example} 令 $f_n(x)=\cos2\pi nx$。则 $f_n\to0$ 于 $L^2$ 的弱拓扑（见\subcref{exercise:ch05-weak-sphere}{item:ch05-orthonormal-weak-zero}），但 $f_n$ 既不几乎处处趋零，也不依测度趋零。
\item\label{item:ch06-spike-weak-example} 令 $f_n(x)=n\chi_{(0,1/n)}(x)$。则 $f_n\to0$ 几乎处处且依测度成立，但对任意 $p$，$f_n$ 都不在 $L^p$ 中弱收敛到零。
\end{enumerate}
\end{exercise}
% SUPPLEMENT: id=v14-ch06-ex22-solution; parent=ch06-ex22; kind=solution+motivation+hypotheses; origin=independent-derivation
\begin{solution}
\textbf{振荡与尖峰的全部核验（重构新增）。}
\begin{enumerate}
\item 对正整数 $n,m$，积化和差公式及整周期积分给
\begin{align*}
\int_0^1\cos(2\pi nx)\cos(2\pi mx)\,\mathrm dx
=\begin{cases}0,&n\ne m,\\ \frac12,&n=m.\end{cases}
\end{align*}
故 $\sqrt2f_n$ 是标准正交序列。对每个 $h\in L^2$，Bessel 不等式给 $\sum_n|\langle h,\sqrt2f_n\rangle|^2\leqslant\|h\|_2^2$，所以这些系数趋零；Hilbert 空间表示定理遂给 $f_n\to0$ 弱收敛。另一方面，每个整周期中 $|\cos t|>1/2$ 的长度占比为 $2/3$，换元并把 $[0,1]$ 分成 $n$ 个周期，得到
\begin{align*}
m\{x\in[0,1]:|f_n(x)|>1/2\}=2/3.
\end{align*}
所以不依测度趋零。若几乎处处趋零，则相应坏集的示性函数几乎处处趋零，且由可积函数一控制，故其积分应趋零，矛盾。因此也不几乎处处趋零。
\item 对 $x=0$，所有 $f_n(x)=0$；对 $x>0$，充分大的 $n$ 满足 $1/n\leqslant x$，也有 $f_n(x)=0$，所以实际上处处趋零。对固定 $\varepsilon>0$，当 $n>\varepsilon$ 时，$m(|f_n|>\varepsilon)=1/n\to0$，故依测度趋零。对每个赋范指数 $1\leqslant p\leqslant\infty$，$f_n\in L^p([0,1])$，且泛函 $T(f)=\int_0^1f$ 连续，因为 $|T(f)|\leqslant\|f\|_p$；可是 $T(f_n)=n\cdot(1/n)=1$。故不弱收敛到零。
\end{enumerate}
\end{solution}
\begin{remark}
\textbf{动机与条件说明（重构新增）。}振荡可使所有固定线性测试趋零，而函数的模长仍在固定比例的区域上很大；尖峰则使坏集越来越小，但保留不消失的积分。第一问的有限总测度把几乎处处收敛提升为依测度收敛，正交性给弱收敛而不需要逐点极限。第二问“任意 $p$”按本节赋范对偶的范围理解为 $1\leqslant p\leqslant\infty$，不能无说明地扩大到 $0<p<1$；该范围外积分一般不是相应拓扑下的连续泛函。
\end{remark}
% END SUPPLEMENT: id=v14-ch06-ex22-solution
% SOURCE p=207; printed=192; object=ch06-ex23; number=23; region=middle-lower
\begin{exercise}\label{exercise:ch06-linfty-semifinite-23}
设 $(X,\mathscr M,\mu)$ 为测度空间。
\begin{definition*}[局部零集]\label{definition:ch06-locally-null}
若 $E\in\mathscr M$，且对每个满足 $\mu(F)<\infty$ 的 $F\in\mathscr M$ 都有 $\mu(E\cap F)=0$，则称 $E$ 为\textbf{局部零集}。
\end{definition*}
% SUPPLEMENT: id=v14-ch06-locally-null-motivation; parent=ch06-ex23; kind=local-definition-motivation+hypotheses; origin=independent-derivation
\begin{remark}
\textbf{局部定义的动机与条件（重构新增）。}有限测度区域是 $L^1$ 函数能够探测的区域；局部零集在每个这种区域内都没有质量。$E,F$ 可测使交集测度有定义；要求对所有有限测度 $F$ 成立，才排除任何局部测试。零集一定局部零，反向在一般测度下不成立，如单点无穷测度空间的全集。这里的“局部”与拓扑邻域无关。
\end{remark}
% END SUPPLEMENT: id=v14-ch06-locally-null-motivation
\begin{definition*}[按局部零集识别的有界函数空间]\label{definition:ch06-local-linfty}
若 $f:X\to\mathbb C$ 可测，定义
\begin{align*}
\|f\|_* =\inf\{a:\{x:|f(x)|>a\}\text{ 是局部零集}\}.
\end{align*}
令 $\mathcal L^\infty=\mathcal L^\infty(X,\mathscr M,\mu)$ 为所有满足 $\|f\|_*<\infty$ 的可测函数构成的空间。若 $\{x:f(x)\ne g(x)\}$ 为局部零集，则把 $f,g\in\mathcal L^\infty$ 视为同一元素。
\end{definition*}
% SUPPLEMENT: id=v14-ch06-local-linfty-motivation; parent=ch06-ex23; kind=local-definition-motivation+hypotheses+source-clarification; origin=independent-derivation
\begin{remark}
\textbf{非负界的约定与动机（重构新增）。}原式把 $a$ 作为模长上界；为涵盖所有测度空间，这里明确在 $a\geqslant0$ 中取下确界，空集合时取 $+\infty$。若允许任意负 $a$，而 $X$ 本身为局部零集，则所有负 $a$ 都会进入原集合，使下确界为 $-\infty$，无法成为范数。采用非负界后，本定义才在这种零空间情形也成立。与通常 $L^\infty$ 相比，此处把所有有限测度测试都看不见的差别识别掉；有限 $\|f\|_*$ 选择具有统一局部本质界的函数，可测性保证超水平集和不等集可测。
\end{remark}
% END SUPPLEMENT: id=v14-ch06-local-linfty-motivation
\begin{enumerate}
\item\label{item:ch06-locally-null-measure} 若 $E$ 为局部零集，则 $\mu(E)$ 为零或无穷。若 $\mu$ 半有限，则每个局部零集都是零集。
\item\label{item:ch06-local-linfty-banach} $\|\cdot\|_*$ 是 $\mathcal L^\infty$ 上的范数，使它成为 Banach 空间。若 $\mu$ 半有限，则 $\mathcal L^\infty=L^\infty$。
\end{enumerate}
\end{exercise}
% SUPPLEMENT: id=v14-ch06-ex23-solution; parent=ch06-ex23; kind=solution+motivation+hypotheses; origin=independent-derivation
\begin{solution}
\textbf{局部零集的商与完备性（重构新增）。}
\begin{enumerate}
\item 若 $\mu(E)<\infty$，在定义中取 $F=E$，便得 $\mu(E)=0$，所以不存在正有限测度的局部零集。若 $\mu$ 半有限而局部零集 $E$ 的测度为无穷，则存在可测 $F\subset E$，$0<\mu(F)<\infty$，但局部零性要求 $\mu(E\cap F)=0$，矛盾。故半有限时局部零集必为零集。
\item 先核验局部零集的运算：可测子集仍局部零；若 $N_j$ 局部零，则对每个有限测度 $F$，
\begin{align*}
\mu\left(F\cap\bigcup_jN_j\right)\leqslant\sum_j\mu(F\cap N_j)=0.
\end{align*}
所以可数并仍局部零。这证明题中识别是等价关系，也使加法和标量乘法与代表无关。若 $f=g$ 在局部零集 $N$ 外成立，则对每个 $a\geqslant0$，两者的超水平集之对称差包含于 $N$；因此是否局部零相同，$\|f\|_* =\|g\|_*$。

记 $M=\|f\|_*<\infty$。可取允许的非负界 $a_j<M+1/j$。于是 $\{|f|>M+1/j\}$ 局部零，而
\begin{align*}
\{|f|>M\}=\bigcup_{j\geqslant1}\{|f|>M+1/j\}
\end{align*}
也局部零。因此每个类都有处处满足 $|f|\leqslant M$ 的有界可测代表：在该超水平集上改为零即可。特别地，$\|f\|_*=0$ 当且仅当 $\{f\ne0\}$ 局部零，恰当它为商空间的零元素。$c\ne0$ 时由超水平集缩放得 $\|cf\|_*=|c|\|f\|_*$，$c=0$ 时直接成立。选 $f,g$ 的上述代表，则处处 $|f+g|\leqslant\|f\|_*+\|g\|_*$，所以 $\|f+g\|_*\leqslant\|f\|_*+\|g\|_*$，也证明此空间线性封闭。故得到范数。

设 $(f_n)$ 关于该范数为 Cauchy 列，先各取一个处处有界可测代表。对每个 $k\geqslant1$，选 $N_k$，使 $m,n\geqslant N_k$ 时 $\|f_m-f_n\|_*\leqslant2^{-k}$。由刚才的界取得性，集合 $\{|f_m-f_n|>2^{-k}\}$ 局部零。将所有 $k$ 及 $m,n\geqslant N_k$ 的这些集合取可数并，得一个局部零集 $N$。在 $N^c$ 上，$(f_n)$ 一致 Cauchy，所以逐点有有限复数极限 $f$，并且一致收敛。该极限有界：在 $N^c$ 上，它与固定的有界函数 $f_{N_1}$ 的差至多为 $1$。在 $N$ 上令 $f=0$；这是可测有界函数。对 $n\geqslant N_k$，令 $m\to\infty$ 得 $|f_n-f|\leqslant2^{-k}$ 于 $N^c$，故 $\|f_n-f\|_*\leqslant2^{-k}$。因此全序列范数收敛，空间完备。

半有限时，由第一问，局部零集恰为零集，所以函数识别、允许的本质界和范数全部与通常 $L^\infty$ 一致，得到等距相同的空间。
\end{enumerate}
\end{solution}
\begin{remark}
\textbf{结论的动机与条件说明（重构新增）。}普通 $L^\infty$ 的证明只需要例外集组成可数并封闭的理想，因此可把零集替换为局部零集。必须在整个证明中维持这种区分：一般情况下，统一例外集 $N$ 未必是 $\mu$ 零集，却已足以控制商范数。有限局部本质界保证有界代表；Cauchy 性保证同一个例外集外的一致极限。半有限性只用于最后与通常 $L^\infty$ 的识别，不是这个新空间完备所需的条件。
\end{remark}
% END SUPPLEMENT: id=v14-ch06-ex23-solution
% SOURCE p=207; printed=192; object=ch06-ex24; number=24; region=lower
\begin{exercise}\label{exercise:ch06-linfty-semifinite-24}
若 $g\in\mathcal L^\infty$（见\cref{exercise:ch06-linfty-semifinite-23}），则
\begin{align*}
\|g\|_* =\sup\left\{\left|\int fg\right|:\|f\|_1=1\right\},
\end{align*}
所以映射 $g\mapsto\phi_g$ 是从 $\mathcal L^\infty$ 到 $(L^1)^*$ 的等距映射。反过来，若如\cref{theorem:ch06-dual-test-criterion}中那样有 $M_\infty(g)<\infty$，则 $g\in\mathcal L^\infty$，且 $M_\infty(g)=\|g\|_*$。
\end{exercise}
% SUPPLEMENT: id=v14-ch06-ex24-solution; parent=ch06-ex24; kind=solution+motivation+hypotheses; origin=independent-derivation
\begin{solution}
\textbf{代表无关性、等距及反向判别（重构新增）。}每个 $f\in L^1$ 的非零集都是有限测度水平集的可数并。若 $N$ 局部零，则其与每个水平集的交为零集，故 $N\cap\{f\ne0\}$ 是 $\mu$ 零集。因此，改变 $g$ 在局部零集上的值不改变 $fg$ 的几乎处处类；改变 $f$ 在零集上的值也不改变乘积类。取 $g$ 的处处有界代表，$|g|\leqslant\|g\|_*$，于是 $fg\in L^1$，且
\begin{align*}
|\phi_g(f)|\leqslant\|g\|_*\|f\|_1.
\end{align*}
线性和连续性随即成立。

令 $M=\|g\|_*$。若 $M>0$，给定 $0<\varepsilon<M$，集合 $A=\{|g|>M-\varepsilon\}$ 不局部零，否则 $M$ 不是下确界。因此存在有限测度可测集 $F$，使 $B=A\cap F$ 满足 $0<\mu(B)<\infty$。取 $f=\mu(B)^{-1}\chi_B\overline{\operatorname{sgn}g}$，则 $\|f\|_1=1$，且
\begin{align*}
\phi_g(f)=\frac1{\mu(B)}\int_B|g|\geqslant M-\varepsilon.
\end{align*}
令 $\varepsilon\downarrow0$，结合上界得 $\|\phi_g\|=M$。$M=0$ 时直接等号成立；空单位球面使用正文的非负上确界约定，等价于在闭单位球取上确界。对 $g-h$ 应用等式，得到线性映射 $g\mapsto\phi_g$ 为等距单射。

再设可测 $g$ 满足原测试假设：所有 $u\in\Sigma$ 均有 $ug\in L^1$，且 $M=M_\infty(g)<\infty$。固定 $\varepsilon>0$，令 $A_\varepsilon=\{|g|>M+\varepsilon\}$。若它不局部零，取有限测度 $F$ 使 $B=A_\varepsilon\cap F$ 正有限测度，再取同一相位测试 $u=\mu(B)^{-1}\chi_B\overline{\operatorname{sgn}g}$。用受 $|u|$ 控制的简单逼近及可积控制 $\mu(B)^{-1}|g|\chi_B$，由 $M$ 的定义得 $|\int ug|\leqslant M$；但直接积分得 $\int ug\geqslant M+\varepsilon$，矛盾。因此每个 $A_\varepsilon$ 局部零，可数并给 $\{|g|>M\}$ 局部零。所以 $g\in\mathcal L^\infty$ 且 $\|g\|_*\leqslant M$。前半部的积分上界用于 $\Sigma$ 又给 $M\leqslant\|g\|_*$，证明相等。
\end{solution}
\begin{remark}
\textbf{动机与条件说明（重构新增）。}改用局部零等价类后，“范数正”的区域按定义就能被某个有限测度集探测，因此等距性不再需要半有限性。$f\in L^1$ 的可数有限测度支集保证积分不受局部零集改变影响，这是配对能下降到商空间的关键。反向判别仍需每个简单乘积可积，才能合法使用控制收敛；有限的统一测试上界把所有超界区域压成局部零集。零空间已包含在单位闭球的公式中。
\end{remark}
% END SUPPLEMENT: id=v14-ch06-ex24-solution
% SOURCE p=207; printed=192; object=ch06-ex25; number=25; region=bottom
\begin{exercise}\label{exercise:ch06-linfty-semifinite-25}
设 $\mu$ 可分解（见\cref{exercise:ch03-decomposable-measure}）。则每个 $\phi\in(L^1)^*$ 都可写为 $\phi(f)=\int fg$，其中 $g\in\mathcal L^\infty$，因而 $(L^1)^*\cong\mathcal L^\infty$（见\cref{exercise:ch06-linfty-semifinite-23,exercise:ch06-linfty-semifinite-24}）。（若 $\mathscr F$ 为 $\mu$ 的一个分解，且 $f\in L^1$，则存在 $\{E_j\}\subset\mathscr F$，使 $f=\sum_{j=1}^\infty f\chi_{E_j}$，其中级数在 $L^1$ 中收敛。）
\end{exercise}
% SUPPLEMENT: id=v14-ch06-ex25-solution; parent=ch06-ex25; kind=solution+motivation+hypotheses; origin=independent-derivation
\begin{solution}
\textbf{逐块表示与可数支集（重构新增）。}固定原定义中的分解 $\mathscr F$：各 $F\in\mathscr F$ 可测且有限测度，两两不交并覆盖 $X$；有限测度集 $E$ 满足 $\mu(E)=\sum_{F\in\mathscr F}\mu(E\cap F)$；若任意子集与每个 $F$ 的交都可测，则该子集可测。下述论证只使用这四条，不额外假设半有限性。

对每个 $F\in\mathscr F$，把 $L^1(F)$ 通过补零嵌入 $L^1(X)$。有限测度的对偶定理给 $g_F\in L^\infty(F)$，使
\begin{align*}
\phi(u)=\int_Fug_F\,\mathrm d\mu\quad(u\in L^1(F)),
\qquad \|g_F\|_\infty\leqslant\|\phi\|.
\end{align*}
各取处处满足 $|g_F|\leqslant\|\phi\|$ 的可测代表（零测度块可取零），在每个 $F$ 上定义 $g=g_F$。因各块不交且覆盖 $X$，$g$ 处处定义。对每个开集 $U\subset\mathbb C$，$g^{-1}(U)\cap F=g_F^{-1}(U)$ 可测；由分解的第四条，$g^{-1}(U)$ 可测。因此 $g$ 可测且处处有界，特别地 $g\in\mathcal L^\infty$。

现在证明原提示。给定 $f\in L^1$，令 $A_n=\{|f|>1/n\}$，其测度有限。分解的第三条给
\begin{align*}
\mu(A_n)=\sum_{F\in\mathscr F}\mu(A_n\cap F)<\infty.
\end{align*}
一个有限的非负和中，正项至多可数：对每个 $k$，大于 $1/k$ 的项只能有限，再取可数并。因此对每个 $n$，只有至多可数个 $F$ 满足 $\mu(A_n\cap F)>0$。把这些块对所有 $n$ 合并，得到至多可数族 $\mathscr F_0$，令 $H=\bigcup_{F\in\mathscr F_0}F$。对有限测度集 $A_n\setminus H$ 再用第三条，每个块交的测度都为零，故 $\mu(A_n\setminus H)=0$。取可数并得 $f=0$ 在 $H^c$ 上几乎处处成立。

若 $\mathscr F_0$ 无限，将其列为互异的 $E_1,E_2,\ldots$；若有限，则使用有限和；若为空，则 $f=0$ 于 $L^1$。记 $H_N=\bigcup_{j=1}^NE_j$。由控制收敛，
\begin{align*}
\left\|f-\sum_{j=1}^Nf\chi_{E_j}\right\|_1
=\int_{H_N^c}|f|\longrightarrow0.
\end{align*}
因此连续性及逐块表示给
\begin{align*}
\phi(f)&=\lim_N\sum_{j=1}^N\phi(f\chi_{E_j})
=\lim_N\int_{H_N}fg
=\int_Xfg,
\end{align*}
最后一步由 $|fg|\leqslant\|\phi\||f|\in L^1$ 及 $f=0$ 于 $H^c$ 几乎处处得到。故表示存在。\cref{exercise:ch06-linfty-semifinite-24}给表示元作为 $\mathcal L^\infty$ 等价类的唯一性与范数等式，所以该同构是自然线性等距同构。
\end{solution}
\begin{remark}
\textbf{动机与逐条件说明（重构新增）。}可分解性把空间组织成有限测度块，使有限区域的对偶表示可逐块拼接。第一条负责局部表示，第二条负责无歧义拼接，第三条使每个 $L^1$ 函数只需可数多个块并保证其余部分确为零，第四条保证不可数拼接仍可测。第三条在原定义中只对有限测度 $E$ 要求，所以证明专门用于有限的 $A_n$ 及 $A_n\setminus H$；不能把它无条件用于所有集合。各块上的零集若作任意不可数并，未必为全局零集；本题的表示元按局部零集识别，避免了这种误推。可分解性比“整个空间由可数个有限测度集覆盖”更宽，正好容纳不可数块的情形。
\end{remark}
% END SUPPLEMENT: id=v14-ch06-ex25-solution
\end{document}
